\subsection{集合的补集。空集}

关系 \((x \notin A \text{ 且 } x \in X)\) 由 C51 可知关于 x 是集合化的。

定义3. 设A是集合X的一个子集。X中不属于A的元素所构成的集合，即集合 \(\mathfrak{E}\mathbf{x}\) ( \(x \notin \mathrm{A}\) 并且 \(x \in \mathrm{X}\) ），称为A在X中的补集，记为 \(\complement_{\mathbf{X}}\mathrm{A}\) 要么 \(\mathrm{X} - \mathrm{A}\) （或记作 \(\complement\mathrm{A}\) （若无混淆风险时）。

设A是集合X的一个子集；关系“ \(x \in X\) 并且 \(x \notin A\) ”和 \(x \in \complement_{X} A\) ”是等价的。因此关系“ \(x \in X\) 并且 \(x \notin \complement_{X} A\) ”等价于“ \(x \in X\) 且 ( \(x \notin X\) 要么 \(x \in A\) ）”，从而等价于 \(x \in A\) 。换言之， \(A = \complement_{\mathbf{X}}(\complement_{\mathbf{X}} A)\) 是一个真关系。类似地，可以证明如果 \(B\) 是……的子集 \(X\) 的显式公理中，关系 \(A \subset B\) 并且 \(\complement_{\mathbf{X}} B \subset \complement_{\mathbf{X}} A\) 是等价的。

定理1. 关系 \((\forall x)(x \notin X)\) 在 \(X\) .

关系 \((\forall x)(x \notin X)\) 蕴含 \((\forall Y)(X \subset Y)\) ；根据外延公理，关系 \((\forall x)(x \notin X)\) 因而关于 \(X\) 是单值的。另一方面，关系 \((\forall x)(x \notin C_{Y}Y)\) 为真，这证明 \((\exists X)(\forall x)(x \notin X)\) 为真。

术语 \(\tau_x((\forall x)(x \notin X))\) 与该函数关系相对应的函数符号表示为 \(\phi\) ，称为空集（*）；关系 \((\forall x)(x \notin X)\) ，它等价于 \(X = \phi\) ，读作：“集合 \(X\) 是空的”。我们有定理 \(x \notin \phi\) , \(\phi \subset X\) , \(\complement_X X = \phi\) , \(\complement_X \phi = X\) 。关系 \(X \subset \phi\) 等价于 \(X = \phi\) 。如果 \(R \{x\}\) 是一个关系，关系 \((\forall x)((x \in \phi) \Longrightarrow R \{x\})\) 为真。

注。不存在一个集合，使得每个对象都是它的元素；换言之，“非 \((\exists X)(\forall x)(x \in X)\) ”是一个定理。因为如果存在这样的集合，那么根据C52，每个关系都将是集体化的。但我们已经看到（第4节），关系 \(x \notin x\) 不是集体化的。

